\section{Gluon helicity distribution}\label{sec:helicity}

Following the same procedure, we can define three operators that can be used to calculate the gluon helicity distribution:
\begin{align}
\Delta {O}^{1}_R(z_2, z_1) &= Z_{11}^2 e^{\overline{\delta { m}}\Delta z}  \epsilon_{ij}  F^{ti}(z_2) L(z_2,z_1) F^{tj}(z_1), \non\\
 \Delta{  O}^{2}_R(z_2,z_1) &=  Z_{22}^2 e^{\overline{\delta { m}}\Delta z} \epsilon_{ij}  F^{zi}(z_2) L(z_2,z_1) F^{zj}(z_1),\non\\
\Delta {O}^{3}_R(z_2,z_1) &=  Z_{11}Z_{22} e^{\overline{\delta { m}}\Delta z}  \epsilon_{ij} F^{ti}(z_2) L(z_2,z_1) F^{zj}(z_1),
\end{align}
where $ \epsilon_{ij} $ is the two-dimensional anti-symmetric tensor.
